
Deep reinforcement learning~(RL) algorithms are predominantly evaluated by comparing their relative performance on a large suite of tasks.
Most published results on deep RL benchmarks compare \emph{point estimates} of aggregate performance such as mean and median scores across tasks, ignoring the statistical uncertainty implied by the use of a finite number of training runs.
Beginning with the Arcade Learning Environment~(ALE), the shift towards computationally-demanding benchmarks has led to the practice of evaluating only a small number of runs per task,
exacerbating the statistical uncertainty in point estimates.
In this paper, we argue that reliable evaluation in the few-run deep RL regime cannot ignore the uncertainty in results without running the risk of slowing down progress in the field. 
We illustrate this point using a case study on the Atari 100k benchmark, where we find substantial discrepancies between conclusions drawn from point estimates alone versus a more thorough statistical analysis.
With the aim of increasing the field's confidence in reported results with \emph{a handful of runs}, we advocate for reporting interval estimates of aggregate performance and propose performance profiles to account for the variability in results, as well as present more robust and  efficient aggregate metrics,
such as interquartile mean scores, to achieve small uncertainty in results.
% , including interval estimates and distribution functions.
Using such statistical tools, we scrutinize performance evaluations of existing algorithms on other widely used RL benchmarks including the ALE, Procgen, and the DeepMind Control Suite, again revealing discrepancies in prior comparisons. 
Our findings call for a change in how we evaluate performance in deep RL, for which we present a more rigorous evaluation methodology, accompanied with an open-source library \href{https://github.com/google-research/rliable}{\emph{rliable}}\footnote{\url{https://github.com/google-research/rliable}}, to prevent unreliable results from stagnating the field. %\footnote{\textbf{Outstanding Paper Award}. Website at \url{https://agarwl.github.io/rliable}}.
% before unreliable results stagnate the field.


% \sout{To this end, we call for the realization that performance on a task is a \emph{random variable} and must be treated as such. This perspective enables researchers to employ statistical tools such as interval estimates and cumulative distribution functions to report the variability in results.} 
% \aaron{While it is widely recognized that task performance estimates are an application of statistical estimation, this recognition rarely comes with any statistical sophistication. In this paper, we show that, in the context of deep RL, a more rigorous methodology is needed to evaluate and compare algorithms}. 